\documentclass[%
  ukrainian, utf8, columnvii, simple,
  hpadding=10mm, vpadding=10mm, reduceheight=2mm,
  floatsection,
]{eskdtext}

\usepackage{eskddstu}
%\usepackage{eskddstu}

\usepackage{fontspec}
\usepackage{xecyr}
\usepackage{xunicode}
\usepackage{xltxtra}
%\usepackage{setspace}

\defaultfontfeatures{Ligatures=TeX}

\setmainfont{Times New Roman}
\newfontfamily\gostfont{opengost-type-b.otf}[AutoFakeSlant=0.26,Scale=MatchLowercase]

\renewcommand{\ESKDfontSetBaseLineStretch}{%
  \renewcommand{\baselinestretch}{\ESKDfontBaseLineStretch}
  \hspace{-1mm}\gostfont}

\linespread{1.5} % 1.5 spacing everywhere
% \setstretch{1.5}   % 1.5 spacing where it is sensible (e.g., not in footnotes)
\usepackage[ukrainian]{babel}

%%% Основні відомості для заповнення%%%

\newcommand{\theWordDone}{Виконав}               % Виконав vs Виконала

\newcommand{\theStudent}{студент}                % студент vs студентка
\newcommand{\theSYear}{4}                        % Курс: 4
\newcommand{\theGroup}{ІО-25}                    % Група
\newcommand{\theStampCode}{ІАЛЦ.467200}          % Код у штампі: код КПІ та код пристрою 



\newcommand{\thesisAuthor}{Льоскін Іван Вадимович}             % Автор дипломного проекту
\newcommand{\thesisAuthorTask}{Льоскіну Івану Вадимовичу}       % Кому видане завдання для дипломного проектування
\newcommand{\thesisAuthorInitials}{Льоскін~І.~В.}              % Прізвище та ініціали автора дипломного проекту

\newcommand{\taskDate}{10~травня 2026 року}      % Дата видачі завдання
\newcommand{\completeDate}{5~червня \the\year\ року}  % Дата здачі закінченого проекту
\newcommand{\approvalDate}{5~червня 2026 року}  % Дата затвердення наказом по університету
\newcommand{\approvalNum}{1212-­с}   % Номер наказу затвердення теми

\newcommand{\thisisHD}           % Завідувач кафедри:  ім’я, прізвище
{\todo{Артем ВОЛОКИТА}}
\newcommand{\thisisHDInitials}   % Завідувач кафедри:  прізвище, ініціали
{Волокита~А.~М.}

\newcommand{\thesisTitle}        % Тема дипломного проєкту
{Кросплатформна клієнт-серверна система агрегації та обробки медико-біологічних даних з використанням хмарних технологій}
%{Інтерфейс користувача для класифікації рукописних цифр}

\newcommand{\thesisOsvitnyaPrograma}     % Освітня програма навчання
{\todo{Комп’ютерні системи та мережі}}

\newcommand{\thesisSpecialtyNumber}    % Номер спеціальності
{123}

\newcommand{\thesisSpecialtyTitle}     % Назва спеціальності
{Комп’ютерна інженерія}

\newcommand{\thesisDegree}             % Освiтньо-квалiфiкацiйний рiвень
{Бакалавр}


\newcommand{\thesisOrganization}       % Назва навчального закладу
{Національний технічний університет України}
\newcommand{\thesisOOrganization}
{"Київський політехнічний інститут імені Ігоря Сікорського"}

\newcommand{\thesisInstitute}           % Факультет
{Факультет інформатики та обчислювальної техніки}

\newcommand{\thesisDepartment}          % Кафедра
{Кафедра обчислювальної техніки}


\newcommand{\supervisorFio}            % Керівник проекту: прізвище, ім’я, по батькові
{Кобилюк Андрій Григорович}
\newcommand{\supervisorInitials}       % Керівник проекту:  прізвище, ініціали
{Кобилюк~А.~Г.}
\newcommand{\supervisorRegalia}        % Керівник проекту: посада, науковий ступінь, вчене звання
{\todo{посада, науковий ступінь}}

\newcommand{\visorFio}                 % Рецензент проекту: прізвище, ім’я, по батькові
{Сергієнко П.А.}
\newcommand{\visorRegalia}             % Рецензент проекту: посада, науковий ступінь, вчене звання
{ас. каф. СПСКС, д-р ф.}

\newcommand{\consultantFio}            % Консультант проекту: прізвище, ім’я, по батькові
{Нікольський С.С.}
\newcommand{\consultantRegalia}        % Консультант проекту: посада, науковий ступінь, вчене звання
{ас. каф. ОТ}
\newcommand{\consultantInitials}       % Консультант проекту: прізвище та ініціали
{Нікольський~С.С.}
\newcommand{\consultantPart}       % Консультант проекту: назва розділу
{Нормоконтроль}
                     

%-------- Вихiднi данi до проєкту -------------
\newcommand{\dataInput}
{
технічна документація, теоретичні та статистичні дані, патенти на винахід, наукові публікації.
}

%-------- Змiст розрахунково-пояснювальної записки -------------
\newcommand{\explanatoryNoteContents}
{
опис предметної області і напрямків дослідження, аналіз та характеристика об’єкту досліджень, аналіз існуючих рішень, опис архітектури системи.
}

%-------- Перелік графічного матеріалу -------------
\newcommand{\graphicMaterial}
{
принципова, структурна та  функціональна схеми системи.
}


%-------- Консультанти розділів проєкту -------------
\newcommand{\consultantsTab}
{
\begin{table}[H]
\begin{center}
\begin{tabular}{ |>{\centering}m{0.5cm}|m{5cm}|>{\centering}m{5.5cm}|>{\centering}m{2cm}|>{\centering}m{2cm}|}
\hline
\rowcolor{lightgray!20}
  &  &  & \multicolumn{2}{c |}{Підпис, дата}  \tabularnewline \cline{4-5}
  
\rowcolor{lightgray!20} 
\multirow{-1}{*}{№} & \multicolumn{1}{c|}{\multirow{-1}{*}{\raggedright Розділ}} &\multicolumn{1}{c|}{\multirow{-1}{*}{\raggedright Консультант}}  & завдання видав & завдання прийняв \tabularnewline
 
\rowcolor{lightgray!20}  
&  &   \multicolumn{1}{c|}{\multirow{-2}{*}{\raggedright\scriptsize(прізвище та ініціали)}}  &  &    \tabularnewline
\hline
   
\rownumber & \consultantPart  & \consultantInitials &  &  \TBstrut \tabularnewline  
\hline

%\rownumber & Назва розділу консультацій   & \consultantInitials &  &  \TBstrut \tabularnewline  % Поле для іншого консультанта
%\hline

\end{tabular}
\end{center}
\end{table}
}


%-------- Календарний план -------------
\newcommand{\calendarPlan}
{
\begin{table}[H]
\setcounter{magicrownumbers}{0}
\begin{center}
\begin{tabular}{|>{\centering}m{0.5cm}|>{\raggedright}m{8cm}|>{\centering}m{5cm}|>{\centering}m{2cm}|}
\hline
\rowcolor{lightgray!20} № & Назва етапу виконання дипломного проекту & Термін виконання & Примітка  \tabularnewline
\hline

\rownumber & Затвердження теми проєкту  & 06.04.2026-07.06.2026 &  \TBstrut \tabularnewline
\hline

\rownumber & Вивчення та аналіз завдання  & 02.05.2026-08.05.2026 &  \TBstrut \tabularnewline
\hline

\rownumber & Розробка архітектури та загальної структури системи  & 09.05.2026-15.05.2026 &  \TBstrut \tabularnewline
\hline

\rownumber & Програмна реалізація системи   & 16.05.2026-29.05.2026 &  \TBstrut \tabularnewline
\hline

\rownumber & Оформлення пояснювальної записки  & 30.05.2026-07.06.2026 &  \TBstrut \tabularnewline
\hline

\rownumber & Передзахист  & 06.06.2026 &  \TBstrut \tabularnewline
\hline

\rownumber & Захист  & 20.06.2026 &  \TBstrut \tabularnewline
\hline

%\rownumber &   &  &  \TBstrut \tabularnewline  % поле для нового етапу
%\hline
\end{tabular}
\end{center}
\end{table}
}


{}
\input{common/stamp-vars}
\ESKDauthor{\thesisAuthorInitials}
\ESKDchecker{\supervisorInitials}
\ESKDnormContr{\consultantInitials}
\ESKDapprovedBy{\thisisHDInitials}

\ESKDtitle{{\fontsize{10}{10.5}\selectfont\thesisTitle}}
\ESKDgroup{\small НТУУ ''КПІ ім. Ігоря Сікорського'', ФІОТ \theGroup}

\newcommand{\stampDocIdBase}{\theStampCode.}

\newcommand{\stampDocId}[2]{\ESKDsignature{\stampDocIdBase#1 #2}}
\newcommand{\stampDocName}[1]{\ESKDdocName{{\fontsize{10}{10.5}\selectfont#1}}}

\makeatletter
\setlength{\ESKD@margin@si}{18mm}
\setlength{\ESKD@margin@so}{7mm}
\setlength{\ESKD@margin@t}{5mm}
\setlength{\ESKD@margin@b}{6mm}
\makeatother

\usepackage[nobottomtitles*]{titlesec}
\usepackage{titletoc}

\titleformat{\section}[display]
  {\filcenter\bfseries\fontsize{18pt}{22pt}\selectfont}
  {\MakeUppercase{\chaptername} \thesection}
  {0pt}
  {\MakeUppercase}[]

\titlespacing{\section}{0pt}{1em plus 0.3em minus 0.06em}{1.5em plus 0.15em}
\newcommand{\sectionbreak}{\clearpage}

\titleformat{\subsection}[hang]
  {\bfseries\fontsize{16pt}{19pt}\selectfont}
  {\thesubsection}
  {1em}{}[]

\titlespacing{\subsection}{\parindent}{1em plus 0.3em minus 0.06em}{1em plus 0.1em}

\titleformat{\subsubsection}[hang]
  {\bfseries}
  {\thesubsubsection}
  {1em}{}[]

\titlespacing{\subsubsection}{\parindent}{1em plus 0.3em minus 0.06em}{1em plus 0.1em}

\titleformat{\paragraph}[hang]
  {\bfseries}
  {\thesubsubsection}
  {1em}{}[]

\titlespacing{\paragraph}{\parindent}{1em plus 0.3em minus 0.06em}{1em plus 0.1em}

\setcounter{secnumdepth}{3}

\newenvironment{unnumbered}{%
  \setcounter{secnumdepth}{0}
}{%
  \setcounter{secnumdepth}{3}
}

% --- TOC entry formatting (match ДСТУ/ESKD style, no bold) ---
\contentsmargin{0pt}

\titlecontents{section}[0pt]
  {\addvspace{0.6em}}
  {РОЗДІЛ~\thecontentslabel.\enspace}
  {}
  {\titlerule*[4pt]{.}\thecontentspage}

\titlecontents{subsection}[2em]
  {}
  {\thecontentslabel\enspace}
  {}
  {\titlerule*[4pt]{.}\thecontentspage}

\titlecontents{subsubsection}[4em]
  {}
  {\thecontentslabel\enspace}
  {}
  {\titlerule*[4pt]{.}\thecontentspage}


\usepackage{csquotes}

\usepackage[%
style=gost-numeric,
language=autobib,
autolang=hyphen,
clearlang=true,
sortcites=true,
sorting=none,
]{biblatex}

\DefineBibliographyStrings{ukrainian}{%
  mediaeresource = {електронний ресурс},
}

\DeclareFieldFormat{url}{Режим доступу до ресурсу: \href{#1}{#1}}
\DeclareFieldFormat{doi}{DOI: \href{https://doi.org/#1}{#1}}

\toggletrue{bbx:gostbibliography}
\renewcommand*{\newblockpunct}{\addperiod\addnbspace-\space\bibsentence}
\AtBeginDocument{\def\bibrangedash{-}}

\addbibresource{bibliography.bib}

\usepackage{enumitem}

\renewcommand{\theenumi}{\arabic{enumi}}
\renewcommand{\labelenumi}{\theenumi)}
\setlist{nosep, topsep=0.5em}
\setlist[itemize]{label=-}

\makeatletter\let\theHpoint\relax\makeatother
\usepackage[linktoc=all, hidelinks]{hyperref}
\usepackage{url}

\usepackage{tikz}

\usetikzlibrary{arrows, automata, shapes.geometric}

\tikzstyle{flow:startstop} = [
  rectangle, rounded corners,
  minimum width=3cm, minimum height=1cm,
  text centered, text width=3cm,
  draw=black, fill=red!30
]

\tikzstyle{flow:io} = [
  trapezium, trapezium left angle=70, trapezium right angle=110,
  minimum width=3cm, minimum height=1cm,
  text centered, text width=3cm,
  draw=black, fill=blue!30
]

\tikzstyle{flow:process} = [
  rectangle,
  minimum width=3cm, minimum height=1cm,
  text centered, text width=3cm,
  draw=black, fill=orange!30
]

\tikzstyle{flow:decision} = [
  diamond,
  minimum width=3cm, minimum height=1cm,
  text centered, text width=3cm,
  draw=black, fill=green!30
]

\tikzstyle{flow:arrow} = [thick,->,>=stealth]

\usepackage{caption}
\usepackage[below]{placeins}
\usepackage{flafter}

\counterwithin{figure}{section}
\counterwithin{table}{section}

\DeclareCaptionLabelSeparator{endash}{\mbox{ - }}
\captionsetup{labelsep=endash, justification=centering, font=normalsize}
\captionsetup[figure]{labelsep=endash, aboveskip=8pt, belowskip=4pt}
\captionsetup[table]{labelsep=endash, position=above, aboveskip=4pt, belowskip=8pt, justification=centering}

\usepackage[toc,nopostdot,nogroupskip,numberedsection=nolabel,nonumberlist]{glossaries-extra}

\makeglossaries

\newabbreviation{ER}{ER}{%
  Entity-Relationship (сутність-зв'язок)}
\newabbreviation{OAuth}{OAuth}{%
  Open Authorization (відкритий протокол авторизації)}
\newabbreviation{URL}{URL}{%
  Uniform Resource Locator (уніфікований локатор ресурсу)}
\newabbreviation{API}{API}{%
  Application Programming Interface (прикладний програмний інтерфейс)}
\newabbreviation{БД}{БД}{%
  база даних}
\newabbreviation{ВООЗ}{ВООЗ}{%
  Всесвітня організація охорони здоров'я}
\newabbreviation{HTTP}{HTTP}{%
  HyperText Transfer Protocol (протокол передачі гіпертексту)}
\newabbreviation{HTTPS}{HTTPS}{%
  HyperText Transfer Protocol Secure (захищений протокол передачі гіпертексту)}
\newabbreviation{JWT}{JWT}{%
  JSON Web Token (веб-токен JSON)}
\newabbreviation{LLM}{LLM}{%
  Large Language Model (велика мовна модель)}
\newabbreviation{LPU}{LPU}{%
  Language Processing Unit (процесор обробки мови)}
\newabbreviation{RAG}{RAG}{%
  Retrieval-Augmented Generation (генерація з доповненим пошуком)}
\newabbreviation{REST}{REST}{%
  Representational State Transfer (передача репрезентативного стану)}
\newabbreviation{RLS}{RLS}{%
  Row-Level Security (безпека на рівні рядків)}
\newabbreviation{SDK}{SDK}{%
  Software Development Kit (набір інструментів розробника)}
\newabbreviation{SQL}{SQL}{%
  Structured Query Language (мова структурованих запитів)}
\newabbreviation{СУБД}{СУБД}{%
  система управління базами даних}
\newabbreviation{UI}{UI}{%
  User Interface (інтерфейс користувача)}
\newabbreviation{ШІ}{ШІ}{%
  штучний інтелект}

\glsfindwidesttoplevelname
\setglossarystyle{long}
\setlength{\glsdescwidth}{.82\textwidth}
\setabbreviationstyle{short-nolong}

\usepackage{longtable}
\usepackage{tabu}

\captionsetup[table]{position=t,justification=centering,slc=off}
\renewcommand{\arraystretch}{1.4}

\usepackage[linesnumbered,lined,boxed]{algorithm2e}


\hyphenpenalty=10000
\exhyphenpenalty=10000
\sloppy

\setlength{\parindent}{1.2cm}


\stampDocName{Опис програми}
\stampDocId{004}{ОП}

\setmonofont{Times New Roman}

\usepackage{verbatim}

\newcommand{\jstpFile}[1]{%
  \textbf{#1}

  \verbatiminput{jstp-src/#1}
}

\begin{document}

\maketitle

\jstpFile{jstp.js}

% \jstpFile{lib/applications.js}

% \jstpFile{lib/common.js}

\jstpFile{lib/connection.js}

% \jstpFile{lib/errors.js}

% \jstpFile{lib/internal-constants.js}

\jstpFile{lib/net.js}

% \jstpFile{lib/remote-proxy.js}

% \jstpFile{lib/serde-fallback.js}

% \jstpFile{lib/serde.js}

\jstpFile{lib/server.js}

% \jstpFile{lib/session.js}

% \jstpFile{lib/simple-connect-policy.js}

% \jstpFile{lib/simple-session-storage-provider.js}

\jstpFile{lib/socket.js}

\jstpFile{lib/stringify.js}

% \jstpFile{lib/tls.js}

\jstpFile{lib/transport-common.js}

% \jstpFile{lib/ws-browser.js}

% \jstpFile{lib/ws-internal.js}

% \jstpFile{lib/ws.js}

% \jstpFile{lib/wss.js}

\textbf{src/common.h}

\begin{verbatim}
#ifndef SRC_COMMON_H_
#define SRC_COMMON_H_

#define THROW_EXCEPTION(ex_type, ex_msg)                                       \
  isolate->ThrowException(v8::Exception::ex_type(                              \
      v8::String::NewFromUtf8(isolate, ex_msg)))

#endif // SRC_COMMON_H_
\end{verbatim}

\textbf{src/message\_parser.h}

\begin{verbatim}
#ifndef SRC_MESSAGE_PARSER_H_
#define SRC_MESSAGE_PARSER_H_

#include <v8.h>

namespace jstp {

namespace message_parser {

// Efficiently parses JSTP messages for transports that require message
// delimiters eliminating the need to split the stream data into parts before
// parsing and allowing to do that in one pass.
v8::Local<v8::String> ParseNetworkMessages(v8::Isolate* isolate,
    const v8::String::Utf8Value& in, v8::Local<v8::Array> out);

}  // namespace message_parser

}  // namespace jstp

#endif // SRC_MESSAGE_PARSER_H_
\end{verbatim}

\textbf{src/message\_parser.cc}

\begin{verbatim}
#include "message_parser.h"

#include <cstddef>
#include <cstring>

#include <v8.h>

#include "common.h"
#include "parser.h"

using std::size_t;
using std::strlen;

using v8::Array;
using v8::Isolate;
using v8::Local;
using v8::String;

using jstp::parser::internal::ParseObject;
using jstp::parser::internal::SkipToNextToken;

namespace jstp {

namespace message_parser {

Local<String> ParseNetworkMessages(Isolate* isolate,
                                  const String::Utf8Value& in,
                                  Local<Array> out) {
  const size_t total_size = in.length();
  size_t parsed_size = 0;
  const char* source = *in;
  const char* curr_chunk = source;
  int index = 0;

  while (parsed_size < total_size) {
    auto chunk_size = strlen(curr_chunk);
    parsed_size += chunk_size + 1;

    if (parsed_size <= total_size) {
      size_t skipped_size = SkipToNextToken(curr_chunk,
          curr_chunk + chunk_size);
      size_t parsed_chunk_size;
      auto result = ParseObject(isolate, curr_chunk + skipped_size,
          curr_chunk + chunk_size, &parsed_chunk_size);

      if (result.IsEmpty()) {
        return Local<String>();
      }

      parsed_chunk_size += skipped_size;
      parsed_chunk_size += SkipToNextToken(curr_chunk + parsed_chunk_size,
          curr_chunk + chunk_size);

      if (parsed_chunk_size != chunk_size) {
        THROW_EXCEPTION(SyntaxError, "Invalid format");
        return Local<String>();
      }

      out->Set(index++, result.ToLocalChecked());
      curr_chunk += chunk_size + 1;
    }
  }

  auto rest = String::NewFromUtf8(isolate, curr_chunk);
  return rest;
}

}  // namespace message_parser

} // namespace jstp
\end{verbatim}

\textbf{src/parser.h}

\begin{verbatim}
#ifndef SRC_PARSER_H_
#define SRC_PARSER_H_

#include <cstddef>

#include <v8.h>

namespace jstp {

namespace parser {

// Deserializes a UTF-8 encoded string into a JavaScript value
// and returns a handle to it.
v8::Local<v8::Value> Parse(v8::Isolate* isolate,
                           const v8::String::Utf8Value& in);

namespace internal {

// Returns count of bytes needed to skip to next token.
size_t SkipToNextToken(const char* str, const char* end);

// Parses an undefined value from `begin` but never past `end` and returns the
// parsed JavaScript value. The `size` is incremented by the number of
// characters the function has used in the string so that the calling side
// knows where to continue from.
v8::MaybeLocal<v8::Value> ParseUndefined(v8::Isolate* isolate,
                                         const char*  begin,
                                         const char*  end,
                                         std::size_t* size);

// Parses a null value from `begin` but never past `end` and returns the parsed
// JavaScript value. The `size` is incremented by the number of characters the
// function has used in the string so that the calling side knows where to
// continue from.
v8::MaybeLocal<v8::Value> ParseNull(v8::Isolate* isolate,
                                    const char*  begin,
                                    const char*  end,
                                    std::size_t* size);

// Parses a boolean value from `begin` but never past `end` and returns the
// parsed JavaScript value. The `size` is incremented by the number of
// characters the function has used in the string so that the calling side
// knows where to continue from.
v8::MaybeLocal<v8::Value> ParseBool(v8::Isolate* isolate,
                                    const char*  begin,
                                    const char*  end,
                                    std::size_t* size);

// Parses a numeric value from `begin` but never past `end` and returns the
// parsed JavaScript value. The `size` is incremented by the number of
// characters the function has used in the string so that the calling side
// knows where to continue from.
v8::MaybeLocal<v8::Value> ParseNumber(v8::Isolate* isolate,
                                      const char*  begin,
                                      const char*  end,
                                      std::size_t* size);

// Parses a string value from `begin` but never past `end` and returns the
// parsed JavaScript value. The `size` is incremented by the number of
// characters the function has used in the string so that the calling side
// knows where to continue from.
v8::MaybeLocal<v8::Value> ParseString(v8::Isolate* isolate,
                                      const char*  begin,
                                      const char*  end,
                                      std::size_t* size);

// Parses an array from `begin` but never past `end` and returns the parsed
// JavaScript value. The `size` is incremented by the number of characters the
// function has used in the string so that the calling side knows where to
// continue from.
v8::MaybeLocal<v8::Value> ParseArray(v8::Isolate* isolate,
                                     const char*  begin,
                                     const char*  end,
                                     std::size_t* size);

// Parses an object key from `begin` but never past `end` and returns
// the parsed JavaScript value. The `size` is incremented by the number
// of characters the function has used in the string so that the calling side
// knows where to continue from.
v8::MaybeLocal<v8::String> ParseKeyInObject(v8::Isolate* isolate,
                                            const char*  begin,
                                            const char*  end,
                                            std::size_t* size);

// Parses a value corresponding to key inside object from `begin`
// but never past `end` and returns the parsed JavaScript value.
// The `size` is incremented by the number of characters the function has used
// in the string so that the calling side knows where to continue from.
v8::MaybeLocal<v8::Value> ParseValueInObject(v8::Isolate* isolate,
                                             const char*  begin,
                                             const char*  end,
                                             std::size_t* size);

// Parses an object from `begin` but never past `end` and returns the parsed
// JavaScript value. The `size` is incremented by the number of characters the
// function has used in the string so that the calling side knows where to
// continue from.
v8::MaybeLocal<v8::Value> ParseObject(v8::Isolate* isolate,
                                      const char*  begin,
                                      const char*  end,
                                      std::size_t* size);

// Parses a decimal number, either integer or float.
v8::MaybeLocal<v8::Value> ParseDecimalNumber(v8::Isolate* isolate,
                                             const char*  begin,
                                             const char*  end,
                                             std::size_t* size,
                                             bool         negate_result);

// Parses an integer number in arbitrary base without prefixes.
v8::Local<v8::Value> ParseIntegerNumber(v8::Isolate* isolate,
                                        const char*  begin,
                                        const char*  end,
                                        std::size_t* size,
                                        int          base,
                                        bool         negate_result);

// Parses an integer number, which is too big to be parsed using
// ParseIntegerNumber, in arbitrary base without prefixes.
v8::Local<v8::Value> ParseBigIntegerNumber(v8::Isolate* isolate,
                                           const char*  begin,
                                           const char*  end,
                                           std::size_t* size,
                                           int          base,
                                           bool         negate_result);

}  // namespace internal

}  // namespace parser

}  // namespace jstp

#endif // SRC_PARSER_H_
\end{verbatim}

\textbf{src/node\_bindings.cc}

\begin{verbatim}
#include <node.h>
#include <v8.h>

#include "common.h"
#include "parser.h"
#include "message_parser.h"

using v8::Array;
using v8::FunctionCallbackInfo;
using v8::HandleScope;
using v8::Isolate;
using v8::Local;
using v8::Object;
using v8::String;
using v8::Value;

namespace jstp {

namespace bindings {

void Parse(const FunctionCallbackInfo<Value>& args) {
  Isolate* isolate = args.GetIsolate();

  if (args.Length() != 1) {
    THROW_EXCEPTION(TypeError, "Wrong number of arguments");
    return;
  }
  if (!args[0]->IsString() && !args[0]->IsUint8Array()) {
    THROW_EXCEPTION(TypeError, "Wrong argument type");
    return;
  }

  HandleScope scope(isolate);

  String::Utf8Value str(args[0]->ToString());
  auto result = jstp::parser::Parse(isolate, str);
  args.GetReturnValue().Set(result);
}

void ParseNetworkMessages(const FunctionCallbackInfo<Value>& args) {
  Isolate* isolate = args.GetIsolate();

  if (args.Length() != 2) {
    THROW_EXCEPTION(TypeError, "Wrong number of arguments");
    return;
  }
  if (!args[0]->IsString() || !args[1]->IsArray()) {
    THROW_EXCEPTION(TypeError, "Wrong argument type");
    return;
  }

  HandleScope scope(isolate);

  String::Utf8Value str(args[0]->ToString());
  auto array = args[1].As<Array>();
  auto result = jstp::message_parser::ParseNetworkMessages(isolate, str, array);
  args.GetReturnValue().Set(result);
}

void Init(Local<Object> target) {
  NODE_SET_METHOD(target, "parse", Parse);
  NODE_SET_METHOD(target, "parseNetworkMessages",
                  ParseNetworkMessages);
}

NODE_MODULE(jstp, Init);

}  // namespace bindings

} // namespace jstp
\end{verbatim}

% \textbf{src/parser.cc}
% 
% \begin{verbatim}
% #include "parser.h"
% 
% #include <cctype>
% #include <cerrno>
% #include <cmath>
% #include <cstddef>
% #include <cstdint>
% #include <cstdlib>
% #include <cstring>
% #include <functional>
% #include <vector>
% 
% #include "common.h"
% #include "unicode_utils.h"
% 
% using std::atof;
% using std::function;
% using std::isalnum;
% using std::isalpha;
% using std::isdigit;
% using std::isxdigit;
% using std::memcpy;
% using std::memset;
% using std::ptrdiff_t;
% using std::size_t;
% using std::strncmp;
% using std::strncpy;
% using std::strtol;
% using std::toupper;
% 
% using v8::Array;
% using v8::False;
% using v8::Integer;
% using v8::Isolate;
% using v8::Local;
% using v8::Maybe;
% using v8::MaybeLocal;
% using v8::NewStringType;
% using v8::Null;
% using v8::Number;
% using v8::Object;
% using v8::String;
% using v8::True;
% using v8::Undefined;
% using v8::Value;
% 
% using jstp::unicode_utils::CodePointToUtf8;
% using jstp::unicode_utils::IsWhiteSpaceCharacter;
% using jstp::unicode_utils::IsLineTerminatorSequence;
% using jstp::unicode_utils::Utf8ToCodePoint;
% using jstp::unicode_utils::IsIdStartCodePoint;
% using jstp::unicode_utils::IsIdPartCodePoint;
% 
% namespace jstp {
% 
% namespace parser {
% 
% // Enumeration of supported JavaScript types used for deserialization
% // function selection.
% enum Type {
%   kUndefined = 0, kNull, kBool, kNumber, kString, kArray, kObject, kDate
% };
% 
% // Parses the type of the serialized JavaScript value at the position `begin`
% // and before `end`. Returns true if it was able to detect the type, false
% // otherwise.
% static bool GetType(const char* begin, const char* end, Type* type);
% 
% // The table of parsing functions indexed with the values of the Type
% // enumeration.
% static constexpr MaybeLocal<Value> (*kParseFunctions[])(Isolate*,
%                                                         const char*,
%                                                         const char*,
%                                                         size_t*) = {
%   &internal::ParseUndefined,
%   &internal::ParseNull,
%   &internal::ParseBool,
%   &internal::ParseNumber,
%   &internal::ParseString,
%   &internal::ParseArray,
%   &internal::ParseObject
% };
% 
% Local<Value> Parse(Isolate* isolate, const String::Utf8Value& in) {
%   const char* str = *in;
%   const size_t length = in.length();
%   const char* end = str + length;
% 
%   Type type;
% 
%   size_t start_pos = internal::SkipToNextToken(str, end);
%   if (!GetType(str + start_pos, end, &type)) {
%     THROW_EXCEPTION(TypeError, "Invalid type");
%     return Undefined(isolate);
%   }
% 
%   size_t parsed_size = 0;
%   MaybeLocal<Value> result = kParseFunctions[type](isolate,
%                                                    str + start_pos,
%                                                    end,
%                                                    &parsed_size);
% 
%   if (result.IsEmpty()) {
%     return Undefined(isolate);
%   }
% 
%   parsed_size += internal::SkipToNextToken(str + start_pos + parsed_size, end);
%   parsed_size += start_pos;
% 
%   if (length != parsed_size) {
%     THROW_EXCEPTION(SyntaxError, "Invalid format");
%     return Undefined(isolate);
%   }
% 
%   return result.ToLocalChecked();
% }
% 
% static bool GetType(const char* begin, const char* end, Type* type) {
%   bool result = true;
%   switch (*begin) {
%     case ',':
%     case ']': {
%       *type = Type::kUndefined;
%       break;
%     }
%     case '{': {
%       *type = Type::kObject;
%       break;
%     }
%     case '[': {
%       *type = Type::kArray;
%       break;
%     }
%     case '\"':
%     case '\'': {
%       *type = Type::kString;
%       break;
%     }
%     case 't':
%     case 'f': {
%       *type = Type::kBool;
%       break;
%     }
%     case 'n': {
%       *type = Type::kNull;
%       if (begin + 4 <= end) {
%         result = (strncmp(begin, "null", 4) == 0);
%       }
%       break;
%     }
%     case 'u': {
%       *type = Type::kUndefined;
%       if (begin + 9 <= end) {
%         result = (strncmp(begin, "undefined", 9) == 0);
%       }
%       break;
%     }
%     case 'N':
%     case 'I': {
%       *type = Type::kNumber;
%       break;
%     }
%     default: {
%       result = false;
%       if (isdigit(*begin) || *begin == '.' || *begin == '+' || *begin == '-') {
%         *type = Type::kNumber;
%         result = true;
%       }
%     }
%   }
%   return result;
% }
% 
% namespace internal {
% 
% // Returns true if `str` points to a multiline comment ending, false otherwise.
% bool IsMultilineCommentEnd(const char* str, size_t* size) {
%   if (str[0] == '*' && str[1] == '/') {
%     *size = 2;
%     return true;
%   }
%   return false;
% }
% 
% // Returns count of bytes needed to skip to current comment ending.
% size_t SkipToCommentEnd(const char* str, const char* end) {
%   bool is_single_line;
% 
%   switch (str[1]) {
%     case '/': {
%       is_single_line = true;
%       break;
%     }
%     case '*': {
%       is_single_line = false;
%       break;
%     }
%     default: {  // In case it is not a comment start
%       return 0;
%     }
%   }
% 
%   function<bool(const char*, size_t*)> end_check_func;
% 
%   if (is_single_line) {
%     end_check_func = IsLineTerminatorSequence;
%   } else {
%     end_check_func = IsMultilineCommentEnd;
%   }
% 
%   const size_t size = end - str;
% 
%   size_t pos = 2;
%   size_t current_size;
% 
%   for (; pos < size; pos++) {
%     if (end_check_func(str + pos, &current_size)) {
%       return pos + current_size;
%     }
%   }
% 
%   if (is_single_line) {
%     return pos;
%   } else {
%     return 0;
%   }
% }
% 
% size_t SkipToNextToken(const char* str, const char* end) {
%   size_t pos = 0;
%   size_t current_size;
%   const size_t size = end - str;
% 
%   while (pos < size) {
%     if (IsWhiteSpaceCharacter(str + pos, &current_size) ||
%         IsLineTerminatorSequence(str + pos, &current_size)) {
%       pos += current_size;
%     } else if (str[pos] == '/') {
%       size_t to_skip = SkipToCommentEnd(str + pos, end);
%       if (!to_skip) {
%         break;
%       }
%       pos += to_skip;
%     } else {
%       break;
%     }
%   }
% 
%   return pos;
% }
% 
% MaybeLocal<Value> ParseUndefined(Isolate*    isolate,
%                                  const char* begin,
%                                  const char* end,
%                                  size_t*     size) {
%   if (*begin == ',' || *begin == ']') {
%     *size = 0;
%   } else if (*begin == 'u') {
%     *size = 9;
%   } else {
%     THROW_EXCEPTION(TypeError, "Invalid format of undefined value");
%     return MaybeLocal<Value>();
%   }
%   return Undefined(isolate);
% }
% 
% MaybeLocal<Value> ParseNull(Isolate*    isolate,
%                             const char* begin,
%                             const char* end,
%                             size_t*     size) {
%   *size = 4;
%   return Null(isolate);
% }
% 
% MaybeLocal<Value> ParseBool(Isolate*    isolate,
%                             const char* begin,
%                             const char* end,
%                             size_t*     size) {
%   MaybeLocal<Value> result;
%   if (begin + 4 <= end && strncmp(begin, "true", 4) == 0) {
%     result = True(isolate);
%     *size = 4;
%   } else if (begin + 5 <= end && strncmp(begin, "false", 5) == 0) {
%     result = False(isolate);
%     *size = 5;
%   } else {
%     THROW_EXCEPTION(TypeError, "Invalid format: expected boolean");
%   }
%   return result;
% }
% 
% MaybeLocal<Value> ParseNumber(Isolate*    isolate,
%                               const char* begin,
%                               const char* end,
%                               size_t*     size) {
%   bool negate_result = false;
%   const char* number_start = begin;
% 
%   if (*begin == '+' || *begin == '-') {
%     negate_result = *begin == '-';
%     number_start++;
%   }
% 
%   int base = 10;
% 
%   if (*number_start == '0') {
%     number_start++;
% 
%     if (*number_start == 'b' || *number_start == 'B') {
%       base = 2;
%       number_start++;
%     } else if (*number_start == 'o' || *number_start == 'O') {
%       base = 8;
%       number_start++;
%     } else if (*number_start == 'x' || *number_start == 'X') {
%       base = 16;
%       number_start++;
%     } else if (isdigit(*number_start)) {
%       THROW_EXCEPTION(SyntaxError,
%           "Legacy octal and non-octal integer literals are not supported");
%       return MaybeLocal<Value>();
%     } else {
%       number_start--;
%     }
%   }
% 
%   MaybeLocal<Value> result;
% 
%   if (base == 10) {
%     result = ParseDecimalNumber(isolate, number_start, end, size,
%                                 negate_result);
%   } else {
%     result = ParseIntegerNumber(isolate, number_start, end, size,
%                                 base, negate_result);
%     if (*size == 0) {
%       THROW_EXCEPTION(SyntaxError, "Empty number value");
%       return MaybeLocal<Value>();
%     }
%   }
%   *size += number_start - begin;
%   return result;
% }
% 
% MaybeLocal<Value> ParseDecimalNumber(Isolate*    isolate,
%                                      const char* begin,
%                                      const char* end,
%                                      size_t*     size,
%                                      bool        negate_result) {
%   char* number_end;
%   double number = strtod(begin, &number_end);
% 
%   if (negate_result) {
%     number = -number;
%   }
% 
%   // strictly allow only "NaN" and "Infinity"
%   if (std::isnan(number)) {
%     if (strncmp(begin + 1, "aN", 2) != 0) {
%       THROW_EXCEPTION(SyntaxError, "Invalid format: expected NaN");
%       return MaybeLocal<Value>();
%     }
%   } else if (std::isinf(number)) {
%     if (strncmp(begin + 1, "nfinity", 7) != 0) {
%       THROW_EXCEPTION(SyntaxError, "Invalid format: expected Infinity");
%       return MaybeLocal<Value>();
%     }
%   }
% 
%   *size = number_end - begin;
%   return Number::New(isolate, number);
% }
% 
% Local<Value> ParseIntegerNumber(Isolate*    isolate,
%                                 const char* begin,
%                                 const char* end,
%                                 size_t*     size,
%                                 int         base,
%                                 bool        negate_result) {
%   char* number_end;
%   long long value = strtoll(begin, &number_end, base);
%   if (errno == ERANGE) {
%     errno = 0;
%     return ParseBigIntegerNumber(isolate, begin, end, size,
%                                  base, negate_result);
%   }
%   if (negate_result) {
%     value = -value;
%   }
%   *size = static_cast<size_t>(number_end - begin);
%   if (value > INT32_MIN && value < INT32_MAX) {
%     return Integer::New(isolate, value);
%   } else {
%     return Number::New(isolate, value);
%   }
% }
% 
% Local<Value> ParseBigIntegerNumber(Isolate*    isolate,
%                                    const char* begin,
%                                    const char* end,
%                                    size_t*     size,
%                                    int         base,
%                                    bool        negate_result) {
%   *size = end - begin;
%   double result = 0.0;
%   char current_digit;
%   double current_digit_value;
%   char base_digit_count = base > 10 ? 10 : base;
%   char base_alpha_count = base > 10 ? base - 10 : 0;
%   for (size_t i = 0; i < *size; i++) {
%     current_digit = toupper(begin[i]);
%     if ((current_digit < '0' || current_digit >= '0' + base_digit_count) &&
%         (current_digit < 'A' || current_digit >= 'A' + base_alpha_count)) {
%       *size = i;
%       break;
%     }
%     current_digit_value = current_digit <= '9' ? current_digit - '0' :
%                                                  current_digit - 'A' + 10;
%     result *= base;
%     result += current_digit_value;
%   }
%   return Number::New(isolate, negate_result ? -result : result);
% }
% 
% static bool GetControlChar(Isolate*    isolate,
%                            const char* str,
%                            size_t*     res_len,
%                            size_t*     size,
%                            char*       write_to);
% 
% MaybeLocal<Value> ParseString(Isolate*    isolate,
%                               const char* begin,
%                               const char* end,
%                               size_t*     size) {
%   *size = end - begin;
%   char* result = nullptr;
% 
%   enum { kApostrophe = 0, kQMarks} string_mode = (*begin == '\'') ?
%                                                  kApostrophe :
%                                                  kQMarks;
%   bool is_ended = false;
%   size_t res_index = 0;
%   size_t out_offset, in_offset;
% 
%   for (size_t i = 1; i < *size; i++) {
%     if ((string_mode == kQMarks     && begin[i] == '\"') ||
%         (string_mode == kApostrophe && begin[i] == '\'')) {
%       is_ended = true;
%       *size = i + 1;
%       break;
%     }
% 
%     if (begin[i] == '\\') {
%       if (!result) {
%         result = new char[*size + 1];
%         memcpy(result, begin + 1, i - 1);
%       }
%       if (IsLineTerminatorSequence(begin + i + 1, &in_offset)) {
%         i += in_offset;
%       } else {
%         bool ok = GetControlChar(isolate, begin + ++i, &out_offset, &in_offset,
%                                  result + res_index);
%         if (!ok) {
%           delete[] result;
%           return MaybeLocal<Value>();
%         }
%         i += in_offset - 1;
%         res_index += out_offset;
%       }
%     } else if (IsLineTerminatorSequence(begin + i, &in_offset)) {
%       delete[] result;
%       THROW_EXCEPTION(SyntaxError, "Unexpected line end in string");
%       return MaybeLocal<Value>();
%     } else {
%       if (result) {
%         result[res_index] = begin[i];
%       }
%       res_index++;
%     }
%   }
% 
%   if (!is_ended) {
%     delete[] result;
%     THROW_EXCEPTION(SyntaxError, "Error while parsing string");
%     return MaybeLocal<Value>();
%   }
% 
%   Local<String> result_str;
%   if (result) {
%     result_str = String::NewFromUtf8(isolate, result,
%         NewStringType::kNormal, static_cast<int>(res_index)).ToLocalChecked();
%     delete[] result;
%   } else {
%     result_str = String::NewFromUtf8(isolate, begin + 1,
%         NewStringType::kNormal, static_cast<int>(*size - 2)).ToLocalChecked();
%   }
%   return result_str;
% }
% 
% static uint32_t ReadHexNumber(const char* str,
%                               size_t required_len,
%                               bool is_limited,
%                               size_t* len,
%                               bool* ok);
% 
% // Parses a Unicode escape sequence after the '\u' part and returns it's
% // code point value. Supports surrogate pairs. Total size of escape
% // sequence (excluding first '\u') is written in `size`.
% static uint32_t ReadUnicodeEscapeSequence(Isolate* isolate,
%                                           const char* str,
%                                           size_t* size,
%                                           bool* ok) {
%   uint32_t result = 0xFFFD;
% 
%   if (isxdigit(str[0])) {
%     result = ReadHexNumber(str, 4, true, nullptr, ok);
%     if (!*ok) {
%       THROW_EXCEPTION(SyntaxError, "Invalid Unicode escape sequence");
%       return 0xFFFD;
%     }
%     *size = 4;
%   } else if (str[0] == '{') {
%     size_t hex_size;
%     result = ReadHexNumber(str + 1, 0, false, &hex_size, ok);
%     if (!*ok || result > 0x10FFFF) {
%       THROW_EXCEPTION(SyntaxError, "Invalid Unicode escape sequence");
%       return 0xFFFD;
%     }
%     *size = hex_size + 2;
%   } else {
%     THROW_EXCEPTION(SyntaxError, "Expected Unicode escape sequence");
%     *ok = false;
%   }
% 
%   // check for surrogate pair
%   if (0xD800 <= result && result <= 0xDBFF) {
%     size_t low_size;
%     if (str[*size] == '\\' && str[*size + 1] == 'u') {
%       uint32_t low_sur = ReadUnicodeEscapeSequence(isolate,
%                                                    str + *size + 2,
%                                                    &low_size, ok);
%       if (!*ok || !(0xDC00 <= low_sur && low_sur <= 0xDFFF)) {
%         return result;
%       }
%       result = ((result - 0xD800) << 10) + low_sur - 0xDC00 + 0x10000;
%       *size += low_size + 2;
%     }
%   }
% 
%   return result;
% }
% 
% // Parses a part of a JavaScript string representation after the backslash
% // character (i.e., an escape sequence without \) into an unescaped control
% // character and writes it to `write_to`.
% // Returns true if no error occured, false otherwise.
% static bool GetControlChar(Isolate*    isolate,
%                            const char* str,
%                            size_t*     res_len,
%                            size_t*     size,
%                            char*       write_to) {
%   *size = 1;
%   *res_len = 1;
%   bool ok;
%   switch (str[0]) {
%     case 'b': {
%       *write_to = '\b';
%       break;
%     }
%     case 'f': {
%       *write_to = '\f';
%       break;
%     }
%     case 'n': {
%       *write_to = '\n';
%       break;
%     }
%     case 'r': {
%       *write_to = '\r';
%       break;
%     }
%     case 't': {
%       *write_to = '\t';
%       break;
%     }
%     case 'v': {
%       *write_to = '\v';
%       break;
%     }
% 
%     case 'x': {
%       *write_to = static_cast<char>(ReadHexNumber(str + 1, 2, true,
%           nullptr, &ok));
%       if (!ok) {
%         THROW_EXCEPTION(SyntaxError, "Invalid hexadecimal escape sequence");
%         return false;
%       }
%       *size = 3;
%       break;
%     }
% 
%     case 'u': {
%       uint32_t symb_code = ReadUnicodeEscapeSequence(isolate,
%                                                      str + 1,
%                                                      size,
%                                                      &ok);
% 
%       if (!ok) {
%         return false;
%       }
%       CodePointToUtf8(symb_code, res_len, write_to);
%       *size += 1;
%       break;
%     }
% 
%     case '0': {
%       if (isdigit(str[1])) {
%         THROW_EXCEPTION(SyntaxError,
%             "Decimal digits after \\0 are not allowed in strings");
%         return false;
%       }
%       *write_to = 0;
%       break;
%     }
% 
%     default: {
%       if ('0' <= str[0] && str[0] <= '7') {
%         THROW_EXCEPTION(SyntaxError,
%             "Octal escape sequences are not allowed in strings");
%         return false;
%       }
%       *write_to = str[0];
%     }
%   }
% 
%   return true;
% }
% 
% // Parses a hexadecimal number with maximal length of max_len (if is_limited true)
% // into uint32_t. Whether the parsing was successful is determined by the value
% // of `ok`. Resulting size of the value will be outputted in len (if is_limited is
% // false).
% static uint32_t ReadHexNumber(const char* str,
%                               size_t required_len,
%                               bool is_limited,
%                               size_t* len,
%                               bool* ok) {
%   static const int8_t xdigit_table[] = {
%     0, 1, 2, 3, 4, 5, 6, 7, 8, 9, // '0' to '9'
%     -1, -1, -1, -1, -1, -1, -1,   // 0x3A to 0x40
%     10, 11, 12, 13, 14, 15,       // 'A' to 'F'
%     // 'G' to 'Z':
%     -1, -1, -1, -1, -1, -1, -1, -1, -1, -1,
%     -1, -1, -1, -1, -1, -1, -1, -1, -1, -1,
%     -1, -1, -1, -1, -1, -1,       // 0x5B to 0x60
%     10, 11, 12, 13, 14, 15,       // 'a' to 'f'
%   };
% 
%   uint32_t result = 0;
%   uint64_t current_value = 0;
%   size_t current_length = 0;
%   char current_digit;
% 
%   *ok = true;
% 
%   while (isxdigit(str[current_length])) {
%     current_digit = str[current_length];
%     current_length++;
%     current_value *= 16;
%     current_value += xdigit_table[current_digit - '0'];
%     if (current_value > UINT32_MAX) {
%       *ok = false;
%       return result;
%     }
%     result = current_value;
%     if (is_limited && current_length == required_len) {
%       break;
%     }
%   }
% 
%   if (is_limited) {
%     if (current_length < required_len) {
%       *ok = false;
%     }
%   } else {
%     if (current_length == 0) {
%       *ok = false;
%     }
%     *len = current_length;
%   }
% 
%   return result;
% }
% 
% MaybeLocal<String> ParseKeyInObject(Isolate*    isolate,
%                                     const char* begin,
%                                     const char* end,
%                                     size_t*     size) {
%   *size = end - begin;
%   Local<String> result;
%   if (begin[0] == '\'' || begin[0] == '"') {
%     Type current_type;
%     bool valid = GetType(begin, end, &current_type);
%     if (valid && current_type == Type::kString) {
%       size_t offset;
%       MaybeLocal<Value> key = ParseString(isolate, begin, end, &offset);
%       if (key.IsEmpty()) {
%         return MaybeLocal<String>();
%       }
%       result = key.ToLocalChecked().As<String>();
%       *size = offset;
%       return result;
%     } else {
%       THROW_EXCEPTION(SyntaxError,
%           "Invalid format in object: key is invalid string");
%       return MaybeLocal<String>();
%     }
%   } else {
%     size_t current_length = 0;
%     size_t cp_size;
%     uint32_t cp;
%     bool ok;
%     char* fallback = nullptr;
%     size_t fallback_length;
%     bool is_escape = false;
%     while (current_length < *size) {
%       if (begin[current_length] == '\\' &&
%           begin[current_length + 1] == 'u') {
%         cp = ReadUnicodeEscapeSequence(isolate, begin + current_length + 2,
%                                        &cp_size, &ok);
%         if (!ok) {
%           return MaybeLocal<String>();
%         }
%         cp_size += 2;
%         if (!fallback) {
%           fallback = new char[*size + 1];
%           memcpy(fallback, begin, current_length);
%           fallback_length = current_length;
%         }
%         is_escape = true;
%       } else {
%         cp = Utf8ToCodePoint(begin + current_length, &cp_size);
%         is_escape = false;
%       }
%       if (current_length == 0 ? IsIdStartCodePoint(cp) :
%                                 IsIdPartCodePoint(cp)) {
%         if (fallback) {
%           if (!is_escape) {
%             memcpy(fallback + fallback_length, begin + current_length, cp_size);
%             fallback_length += cp_size;
%           } else {
%             size_t fallback_cp_size;
%             CodePointToUtf8(cp, &fallback_cp_size, fallback + fallback_length);
%             fallback_length += fallback_cp_size;
%           }
%         }
%         current_length += cp_size;
%       } else {
%         if (current_length != 0) {
%           if (!fallback) {
%             result = String::NewFromUtf8(isolate, begin,
%                                          NewStringType::kInternalized,
%                                          static_cast<int>(current_length))
%                                              .ToLocalChecked();
%           } else {
%             result = String::NewFromUtf8(isolate, fallback,
%                                          NewStringType::kInternalized,
%                                          static_cast<int>(fallback_length))
%                                              .ToLocalChecked();
%           }
%           break;
%         } else {
%           THROW_EXCEPTION(SyntaxError, "Unexpected identifier");
%           return MaybeLocal<String>();
%         }
%       }
%     }
%     *size = current_length;
%     return result;
%   }
% }
% 
% MaybeLocal<Value> ParseValueInObject(Isolate*    isolate,
%                                      const char* begin,
%                                      const char* end,
%                                      size_t*     size) {
%   Type current_type;
%   bool valid = GetType(begin, end, &current_type);
%   if (valid) {
%     return (kParseFunctions[current_type])(isolate, begin, end, size);
%   } else {
%     THROW_EXCEPTION(TypeError, "Invalid type in object");
%     return MaybeLocal<Value>();
%   }
% }
% 
% MaybeLocal<Value> ParseObject(Isolate*    isolate,
%                               const char* begin,
%                               const char* end,
%                               size_t*     size) {
%   bool key_mode = true;
%   *size = end - begin;
%   MaybeLocal<String> current_key;
%   MaybeLocal<Value> current_numeric_key;
%   MaybeLocal<Value> current_value;
%   size_t current_length = 0;
%   auto result = Object::New(isolate);
%   bool has_ended = false;
% 
%   for (size_t i = 1; i < *size; i++) {
%     if (key_mode) {
%       i += SkipToNextToken(begin + i, end);
%       if (begin[i] == '}') {
%         *size = i + 1;
%         has_ended = true;
%         break;
%       }
%       if (!isdigit(begin[i])) {
%         current_key = ParseKeyInObject(isolate, begin + i, end,
%             &current_length);
%       } else {
%         current_numeric_key = ParseNumber(isolate, begin + i, end,
%             &current_length);
%         current_key = current_numeric_key.IsEmpty() ? MaybeLocal<String>() :
%             current_numeric_key.ToLocalChecked()->ToString(
%                 isolate->GetCurrentContext());
%       }
%       if (current_key.IsEmpty()) {
%         return MaybeLocal<Value>();
%       }
%       i += current_length;
%       i += SkipToNextToken(begin + i, end);
%       if (begin[i] != ':') {
%         THROW_EXCEPTION(SyntaxError, "Unexpected token");
%         return MaybeLocal<Value>();
%       }
%     } else {
%       i += SkipToNextToken(begin + i, end);
%       if (begin[i] == ',') {
%         THROW_EXCEPTION(SyntaxError, "Value is missing in object");
%         return MaybeLocal<Value>();
%       }
%       current_value = ParseValueInObject(isolate,
%                                          begin + i,
%                                          end,
%                                          &current_length);
%       if (current_value.IsEmpty()) {
%         return current_value;
%       }
%       Local<Value> value = current_value.ToLocalChecked();
%       if (!value->IsUndefined()) {
%         Maybe<bool> is_ok = result->Set(isolate->GetCurrentContext(),
%                                         current_key.ToLocalChecked(),
%                                         value);
%         if (is_ok.IsNothing()) {
%           THROW_EXCEPTION(Error, "Cannot add property to object");
%           return MaybeLocal<Value>();
%         }
%       }
%       i += current_length;
%       i += SkipToNextToken(begin + i, end);
%       if (begin[i] != ',' && begin[i] != '}') {
%         THROW_EXCEPTION(SyntaxError, "Invalid format in object");
%         return MaybeLocal<Value>();
%       } else if (begin[i] == '}') {
%         *size = i + 1;
%         has_ended = true;
%         break;
%       }
%     }
%     key_mode = !key_mode;
%   }
% 
%   if (!has_ended) {
%     THROW_EXCEPTION(SyntaxError, "Missing closing brace in object");
%     return MaybeLocal<Value>();
%   }
% 
%   return result;
% }
% 
% MaybeLocal<Value> ParseArray(Isolate*    isolate,
%                              const char* begin,
%                              const char* end,
%                              size_t*     size) {
%   auto array = Array::New(isolate);
%   size_t current_length = 0;
%   *size = end - begin;
% 
%   bool is_empty = true;
%   bool has_ended = false;
% 
%   size_t current_element = 0;
%   Type current_type;
% 
%   for (size_t i = 1; i < *size; i++) {
%     i += SkipToNextToken(begin + i, end);
%     if (is_empty && begin[i] == ']') {  // In case of empty array
%       *size = i + 1;
%       return array;
%     }
% 
%     bool valid = GetType(begin + i, end, &current_type);
%     if (valid) {
%       MaybeLocal<Value> t = kParseFunctions[current_type](isolate,
%                                                           begin + i,
%                                                           end,
%                                                           &current_length);
%       if (t.IsEmpty()) {
%         return t;
%       }
%       if (!(current_type == Type::kUndefined && begin[i] == ']')) {
%         array->Set(static_cast<uint32_t>(current_element++),
%                    t.ToLocalChecked());
%         is_empty = false;
%       }
% 
%       i += current_length;
%       i += SkipToNextToken(begin + i, end);
% 
%       current_length = 0;
% 
%       if (begin[i] != ',' && begin[i] != ']') {
%         THROW_EXCEPTION(SyntaxError, "Invalid format in array: missed comma");
%         return MaybeLocal<Value>();
%       } else if (begin[i] == ']') {
%         *size = i + 1;
%         has_ended = true;
%         break;
%       }
%     } else {
%       THROW_EXCEPTION(TypeError, "Invalid type in array");
%       return MaybeLocal<Value>();
%     }
%   }
% 
%   if (!has_ended) {
%     THROW_EXCEPTION(SyntaxError, "Missing closing bracket in array");
%     return MaybeLocal<Value>();
%   }
% 
%   return array;
% }
% 
% }  // namespace internal
% 
% }  // namespace parser
% 
% } // namespace jstp
% \end{verbatim}

% \textbf{src/unicode\_utils.h}
% 
% \begin{verbatim}
% #ifndef SRC_UNICODE_UTILS_H_
% #define SRC_UNICODE_UTILS_H_
% 
% #include <cstddef>
% #include <cstdint>
% 
% namespace jstp {
% 
% namespace unicode_utils {
% 
% // Returns true if `str` points to a valid Line Terminator Sequence code point,
% // false otherwise. `size` will receive the number of bytes used by this
% // code point (1, 2, 3).
% bool IsLineTerminatorSequence(const char* str, std::size_t* size);
% 
% // Returns true if `str` points to a valid White space code point,
% // false otherwise. `size` will receive the number of bytes used by this
% // code point (1, 2, 3).
% bool IsWhiteSpaceCharacter(const char* str, std::size_t* size);
% 
% // Encodes a Unicode code point in UTF-8 and writes it to `write_to`.
% // `size` will receive the number of bytes used (1, 2, 3 or 4).
% void CodePointToUtf8(unsigned int c, std::size_t* size, char* write_to);
% 
% // Decodes a UTF-8 encoded Unicode code point starting at the `begin`.
% // `size` will receive the number of bytes the code point occupies.
% std::uint32_t Utf8ToCodePoint(const char* begin, std::size_t* size);
% 
% // Checks whether the given Unicode code point is a valid IdentifierStart.
% bool IsIdStartCodePoint(std::uint32_t cp);
% 
% // Checks whether the given Unicode code point is a valid IdentifierPart.
% bool IsIdPartCodePoint(std::uint32_t cp);
% 
% }  // namespace unicode_utils
% 
% }  // namespace jstp
% 
% #endif // SRC_UNICODE_UTILS_H_
% \end{verbatim}
% 
% \textbf{src/unicode\_utils.cc}
% 
% \begin{verbatim}
% #include "unicode_utils.h"
% 
% #include <cstddef>
% #include <cstdint>
% 
% #ifdef _PARSER_USE_FULL_TABLES_
% #include "unicode_tables.h"
% #else
% #include "unicode_range_tables.h"
% #endif
% 
% using std::size_t;
% using std::uint32_t;
% 
% namespace jstp {
% 
% namespace unicode_utils {
% 
% bool IsLineTerminatorSequence(const char* str, size_t* size) {
%   if (str[0] == '\x0D' && str[1] == '\x0A') {
%     *size = 2;
%     return true;
%   } else if (str[0] == '\x0D' || str[0] == '\x0A') {
%     *size = 1;
%     return true;
%   } else if (str[0] == '\xE2' &&
%              str[1] == '\x80' &&
%             (str[2] == '\xA8' ||
%              str[2] == '\xA9')) {
%     *size = 3;
%     return true;
%   }
%   return false;
% }
% 
% bool IsWhiteSpaceCharacter(const char* str, size_t* size) {
%   if (str[0] == '\x09' ||
%       str[0] == '\x0B' ||
%       str[0] == '\x0C' ||
%       str[0] == '\x20' ||
%       str[0] == '\xA0') {
%     *size = 1;
%     return true;
%   } else if (str[0] == '\xC2' && str[1] == '\xA0') {
%     *size = 2;
%     return true;
%   } else {
%     bool is_multibyte_space = false;
%     switch (str[0]) {
%       case '\xE1': {
%         if (str[1] == '\xBB' && str[2] == '\xBF') {
%           is_multibyte_space = true;
%         }
%         break;
%       }
%       case '\xE2': {
%         if ((str[1] == '\x80' &&
%             ((static_cast<unsigned char>(str[2]) & 0x7F) <= 0xA ||
%                                          str[2] == '\xAF'))     ||
%             (str[1] == '\x81' && str[2] == '\x9F')) {
%           is_multibyte_space = true;
%         }
%         break;
%       }
%       case '\xE3': {
%         if (str[1] == '\x80' && str[2] == '\x80') {
%           is_multibyte_space = true;
%         }
%         break;
%       }
%       case '\xEF': {
%         if (str[1] == '\xBB' && str[2] == '\xBF') {
%           is_multibyte_space = true;
%         }
%         break;
%       }
%     }
%     if (is_multibyte_space) {
%       *size = 3;
%       return true;
%     }
%   }
%   return false;
% }
% 
% void CodePointToUtf8(unsigned int c, size_t* size, char* write_to) {
%   char* b = write_to;
%   if (c < 0x80) {
%     *b++ = c;
%     *size = 1;
%   } else if (c < 0x800) {
%     *b++ = 192 + c / 64;
%     *b++ = 128 + c % 64;
%     *size = 2;
%   } else if (c - 0xd800u < 0x800) {
%     CodePointToUtf8(0xFFFD, size, write_to);
%     return;
%   } else if (c < 0x10000) {
%     *b++ = 224 + c / 4096;
%     *b++ = 128 + c / 64 % 64;
%     *b++ = 128 + c % 64;
%     *size = 3;
%   } else if (c < 0x110000) {
%     *b++ = 240 + c / 262144;
%     *b++ = 128 + c / 4096 % 64;
%     *b++ = 128 + c / 64 % 64;
%     *b++ = 128 + c % 64;
%     *size = 4;
%   } else {
%     CodePointToUtf8(0xFFFD, size, write_to);
%     return;
%   }
% }
% 
% uint32_t Utf8ToCodePoint(const char* begin, size_t* size) {
%   auto str = reinterpret_cast<const unsigned char*>(begin);
%   uint32_t result = 0;
%   *size = 1;
%   if (*str < 0x80) {
%     return *str;
%   } else if ((*str & 0xE0) == 0xC0) {
%     *size = 2;
%     result += (*str & 0x1F) << 6;
%   } else if ((*str & 0xF0) == 0xE0) {
%     *size = 3;
%     result += (*str & 0x0F) << 12;
%   } else if ((*str & 0xF8) == 0xF0) {
%     *size = 4;
%     result += (*str & 0x07) << 18;
%   } else {
%     return 0xFFFD;
%   }
%   for (size_t i = 2; i <= *size; i++) {
%     str++;
%     if ((*str & 0xC0) != 0x80) {
%       *size = i;
%       return 0xFFFD;
%     }
%     result += (*str & 0x3F) << ((*size - i) * 6);
%   }
%   return result;
% }
% 
% #ifdef _PARSER_USE_FULL_TABLES_
% 
% bool IsIdStartCodePoint(uint32_t cp) {
%   return ID_START_FULL[cp];
% }
% 
% bool IsIdPartCodePoint(uint32_t cp) {
%   return ID_CONTINUE_FULL[cp];
% }
% 
% #undef _PARSER_USE_FULL_TABLES_
% #else
% 
% bool search_cp(uint32_t cp, const unicode_range* ranges, size_t size) {
%   size_t low = 0;
%   size_t high = size;
%   while (high > low) {
%     size_t mid = low + (high - low) / 2;
%     unicode_range range = ranges[mid];
%     if (range.start <= cp && cp <= range.end) {
%       return true;
%     }
%     if (cp < range.start) {
%       high = mid;
%     } else {
%       low = mid + 1;
%     }
%   }
%   return false;
% }
% 
% bool IsIdStartCodePoint(uint32_t cp) {
%   if (('a' <= cp && cp <= 'z') ||
%       ('A' <= cp && cp <= 'Z') ||
%       cp == '_' ||
%       cp == '$') {
%     return true;
%   }
%   return search_cp(cp, ID_START_RANGES, ID_START_RANGES_COUNT);
% }
% 
% bool IsIdPartCodePoint(uint32_t cp) {
%   if (('a' <= cp && cp <= 'z') ||
%       ('A' <= cp && cp <= 'Z') ||
%       ('0' <= cp && cp <= '9') ||
%       cp == '_' ||
%       cp == '$') {
%     return true;
%   }
%   return search_cp(cp, ID_CONTINUE_RANGES, ID_CONTINUE_RANGES_COUNT) ||
%                    cp == 0x200C || cp == 0x200D; // ZWNJ, ZWJ
% }
% 
% #endif
% 
% 
% }  // namespace unicode_utils
% 
% } // namespace jstp
% \end{verbatim}

\end{document}
